\documentclass{article}

\usepackage{amsmath,amssymb}
\usepackage{xurl}
\usepackage[hidelinks]{hyperref}

% Shared commands for notes in docs/paper-gaps/.

\DeclareUnicodeCharacter{2080}{\ifmmode{}_{0}\else\textsubscript{0}\fi}
\DeclareUnicodeCharacter{2081}{\ifmmode{}_{1}\else\textsubscript{1}\fi}
\DeclareUnicodeCharacter{2082}{\ifmmode{}_{2}\else\textsubscript{2}\fi}
\DeclareUnicodeCharacter{2099}{\ifmmode{}_{n}\else\textsubscript{n}\fi}

\newcommand{\Lean}{\textsc{Lean}}
\ExplSyntaxOn
\NewDocumentCommand{\leanid}{m}{
  \ifmmode
    \textsf{\detokenize{#1}}
  \else
    \group_begin:
    \urlstyle{sf}
    \tl_set:Nn \l_tmpa_tl {#1}
    \tl_replace_all:Nnn \l_tmpa_tl {\_} {_}
    \exp_args:NV \path \l_tmpa_tl
    \group_end:
  \fi
}
\ExplSyntaxOff
\providecommand{\tr}{\operatorname{tr}}
\newcommand{\tnleanIssueBase}{https://github.com/LionSR/TNLean/issues/}
\newcommand{\tnleanPrBase}{https://github.com/LionSR/TNLean/pull/}
\newcommand{\tnleanDocBase}{https://sirui-lu.com/TNLean/docs/}
\newcommand{\ghissue}[1]{\href{\tnleanIssueBase#1}{GitHub issue \##1}}
\newcommand{\ghpr}[1]{\href{\tnleanPrBase#1}{PR \##1}}
\newcommand{\doclink}[1]{\href{\tnleanDocBase#1.html}{\path{#1}}}

% Dirac notation and the identity operator, as in blueprint/src/macros/common.tex.
% \providecommand keeps note-local definitions authoritative.
\usepackage{dsfont}
\providecommand{\ket}[1]{|#1\rangle}
\providecommand{\bra}[1]{\langle #1|}
\providecommand{\braket}[2]{\langle #1 | #2 \rangle}
\providecommand{\ketbra}[2]{|#1\rangle\!\langle #2|}
\providecommand{\Id}{\mathds{1}}
\providecommand{\one}{\mathds{1}}
\providecommand{\C}{\mathbb{C}}
\providecommand{\MN}[1]{M_{#1}(\C)}
\providecommand{\MD}{M_D(\C)}

% External referencing for paper sources.  The build helper writes sanitized
% auxiliary files under build/paper-aux/ and excludes duplicated source labels.
\usepackage{xr}

% Import a generated paper auxiliary file with a caller-supplied label prefix.
% The candidate paths support builds from docs/paper-gaps/, the repository root,
% and build/paper-aux/ itself.
\newcommand{\importpaperaux}[2]{%
  \IfFileExists{../../build/paper-aux/#2.aux}{%
    \externaldocument[#1]{../../build/paper-aux/#2}%
  }{%
    \IfFileExists{build/paper-aux/#2.aux}{%
      \externaldocument[#1]{build/paper-aux/#2}%
    }{%
      \IfFileExists{#2.aux}{%
        \externaldocument[#1]{#2}%
      }{}%
    }%
  }%
}

% General external reference macros
\newcommand{\paperref}[2]{\ref{#1-#2}}
\newcommand{\papereqref}[2]{(\ref{#1-#2})}

% CPSV16 source (1606.00608 / MPDO-22-12-17-2.tex) external document and shortcuts
\importpaperaux{cpsv16-}{1606.00608/MPDO-22-12-17-2-tnlean}
\newcommand{\cpsvref}[1]{\paperref{cpsv16}{#1}}
\newcommand{\cpsveqref}[1]{\papereqref{cpsv16}{#1}}

% Verdict marker: \gapnote{<kind>}{<status>}. Kind is the policy
% classification (clarification, local-correction, scope-restriction,
% unfaithful, false-source, open-gap); status is open, wip (elimination
% underway), resolved, or historical. Machine-read by the site index;
% prints a verdict line.
\newcommand{\gapnote}[2]{\par\noindent\textbf{Verdict:} #1 (#2).\par}


\title{Circle and Nonzero-Complex Coefficients in Degree-Two Cohomology}
\author{}
\date{2026-08-29}

\begin{document}
\maketitle
\gapnote{clarification}{resolved}

\section{Source context}

For a finite on-site symmetry group $G$, Ref.~\cite[lines~1927--1931]{FrancoRubio2025MPUGauging}
labels the fully unbroken phases by $H^2(G,\mathbb C^\times)$, or equivalently
by $H^2(G,U(1))$.  Its footnote regards the two coefficient groups as
topological abelian groups and invokes an induced comparison in sheaf
cohomology over a CW complex.

The formal development instead uses algebraic group cohomology.  It regards
$G$ as finite and discrete and computes cohomology from inhomogeneous cochains
with trivial coefficient action.  The sheaf-cohomology sentence does not
directly identify these algebraic cochain complexes.

\section{Finite-discrete comparison}

Let $i:U(1)\longrightarrow\mathbb C^\times$ be inclusion and let
$q:\mathbb C^\times\longrightarrow U(1)$ be the phase homomorphism
$q(z)=z/|z|$.  Then
\begin{align}
    q\circ i
        & =\operatorname{id}_{U(1)},
    &
    \frac{z}{i(q(z))}
        & =|z|\in\mathbb R_{>0}.
    \label{eq:fbc25_circle_polar_retraction}
\end{align}
Consequently, if $\omega$ is a $\mathbb C^\times$-valued $2$-cocycle, then
$q\circ\omega$ is a $U(1)$-valued $2$-cocycle and its radial residual is
\begin{align}
    \frac{\omega(g,h)}{i(q(\omega(g,h)))}
        & =|\omega(g,h)|.
    \label{eq:fbc25_circle_radial_residual}
\end{align}
The right-hand side of
(\ref{eq:fbc25_circle_radial_residual}) is a positive generalized cocycle on
the one-point block space.  The finite-group positive-cocycle theorem gives a
positive $1$-cochain $\chi$ such that
\begin{align}
    |\omega(g,h)|
        & =\frac{\chi(g)\chi(h)}{\chi(gh)}.
    \label{eq:fbc25_circle_positive_coboundary}
\end{align}
Thus $\omega$ and $i\circ q\circ\omega$ are cohomologous.  Every
$\mathbb C^\times$-valued class therefore has a unit-modulus representative.

Equation~(\ref{eq:fbc25_circle_polar_retraction}) also shows, without assuming
that $G$ is finite, that the map induced by $i$ is injective: the map induced
by $q$ is its left inverse.  Combining this injectivity with the finite-group
surjectivity above gives
\begin{align}
    H^2(G,U(1))
        & \cong H^2(G,\mathbb C^\times)
        \qquad (G\text{ finite}),
    \label{eq:fbc25_circle_h2_equivalence}
\end{align}
for algebraic group cohomology with trivial coefficient actions.

\section{Formal treatment}

The formal proof implements the phase map, the polar identity, the positive
one-block trivialization, and the induced maps on the canonical degree-two
group cohomology object.
\footnote{Formal declarations:
\leanid{TNLean.Algebra.ScalarCocycle.cohomologousTo\_circlePhaseInclusion}
and \leanid{TNLean.Algebra.circleH2EquivComplexUnitsH2} in
\doclink{TNLean/Algebra/CircleCohomology}.}

\textbf{Verdict: resolved clarification.}
The source's finite-group equivalence is available for the algebraic
cohomology used in the classification, but it is established here by a direct
polar-decomposition argument rather than by identifying it with the sheaf
cohomology comparison in the footnote.

\bibliographystyle{alpha}
\bibliography{references}

\end{document}
